% Source: https://artofproblemsolving.com/community/c3h2433694_forester2015s_problems
% Stable identity: aops-topic-2433694-q31
\textbf{Counting directly.} Of each of the $12$ edges of the big cube, only the middle third survives: $12$ edges. At each of the $8$ corners, removing the small cube leaves a notch.

\begin{asy}
usepackage("amsmath");
size(5cm);
defaultpen(fontsize(10pt));
pair dx=(1,0), dy=(0.45,0.3), dz=(0,1);
pair P(real x,real y,real z){return x*dx+y*dy+z*dz;}
// visible faces of the solid: front y=0, right x=3, top z=3 (cross shapes) and the notch faces
path cross(pair o, pair u, pair v){ return o+1u--o+2u--o+2u+1v--o+3u+1v--o+3u+2v--o+2u+2v--o+2u+3v--o+1u+3v--o+1u+2v--o+2v--o+1v--o+1u+1v--cycle; }
// painter's order: the back notches first, then the front notches, then the outer cross-shaped faces
// notch at the back-top-right corner: faces x=2, z=2 (y=2 hidden); partly hidden behind the edge cubes in front of it
filldraw(P(2,2,2)--P(3,2,2)--P(3,3,2)--P(2,3,2)--cycle,rgb(1,0.85,0.6),linewidth(0.8));
filldraw(P(2,2,2)--P(2,3,2)--P(2,3,3)--P(2,2,3)--cycle,rgb(1,0.9,0.7),linewidth(0.8));
// notch at the back-bottom-right corner: face x=2; notch at the back-top-left corner: face z=2
filldraw(P(2,2,0)--P(2,3,0)--P(2,3,1)--P(2,2,1)--cycle,rgb(1,0.9,0.7),linewidth(0.8));
filldraw(P(0,2,2)--P(1,2,2)--P(1,3,2)--P(0,3,2)--cycle,rgb(1,0.85,0.6),linewidth(0.8));
// notch at the front-top-right corner (cube [2,3]x[0,1]x[2,3] removed): faces x=2, y=1, z=2
filldraw(P(2,0,2)--P(3,0,2)--P(3,1,2)--P(2,1,2)--cycle,rgb(1,0.85,0.6),linewidth(0.8));
filldraw(P(2,1,2)--P(3,1,2)--P(3,1,3)--P(2,1,3)--cycle,rgb(1,0.75,0.45),linewidth(0.8));
filldraw(P(2,0,2)--P(2,1,2)--P(2,1,3)--P(2,0,3)--cycle,rgb(1,0.9,0.7),linewidth(0.8));
// notch at the front-top-left corner: faces x=1 (hidden), y=1, z=2
filldraw(P(0,0,2)--P(1,0,2)--P(1,1,2)--P(0,1,2)--cycle,rgb(1,0.85,0.6),linewidth(0.8));
filldraw(P(0,1,2)--P(1,1,2)--P(1,1,3)--P(0,1,3)--cycle,rgb(1,0.75,0.45),linewidth(0.8));
// notch at the front-bottom-right corner: faces x=2, y=1 (z=1 hidden)
filldraw(P(2,1,0)--P(3,1,0)--P(3,1,1)--P(2,1,1)--cycle,rgb(1,0.75,0.45),linewidth(0.8));
filldraw(P(2,0,0)--P(2,1,0)--P(2,1,1)--P(2,0,1)--cycle,rgb(1,0.9,0.7),linewidth(0.8));
// notch at the front-bottom-left corner: face y=1 (x=1 and z=1 hidden)
filldraw(P(0,1,0)--P(1,1,0)--P(1,1,1)--P(0,1,1)--cycle,rgb(1,0.75,0.45),linewidth(0.8));
filldraw(cross(P(0,0,0),dx,dz),rgb(0.9,0.9,1),linewidth(0.8));
filldraw(cross(P(3,0,0),dy,dz),rgb(0.8,0.85,1),linewidth(0.8));
filldraw(cross(P(0,0,3),dx,dy),rgb(0.95,0.95,1),linewidth(0.8));
// highlight the 9 edges created at the front-top-right corner
pen hp=red+linewidth(1.4);
draw(P(2,1,2)--P(3,1,2),hp); draw(P(2,1,2)--P(2,1,3),hp); draw(P(2,1,2)--P(2,0,2),hp);
draw(P(3,0,2)--P(3,1,2)--P(3,1,3),hp); draw(P(2,0,3)--P(2,1,3)--P(3,1,3),hp); draw(P(2,0,2)--P(3,0,2),hp); draw(P(2,0,2)--P(2,0,3),hp);
\end{asy}

The notch has three new square faces meeting at the inner corner, which gives $3$ edges there, and on each of the three old faces around the corner the boundary of the notch consists of $2$ unit segments, which gives $6$ more edges ($9$ edges per corner, shown in red for one corner). So the solid has
\[12+8\cdot9=84\]
edges. The answer is $\boxed{\textbf{(D)}\ 84}$.

\textbf{Check with Euler's formula.} Each corner contributes $7$ vertices (the inner corner, $3$ points on old edges and $3$ points on old faces), so $V=56$; the faces are the $6$ old faces and $3$ new ones per corner, $F=30$. Indeed $V-E+F=56-84+30=2$.
